\textbf{Phần 1. Trắc nghiệm lựa chọn}
\setcounterpageref{ex}{0}

\begin{ex} %[5D10V56]
Bác Hồ đọc bản Tuyên ngôn Độc lập khai sinh ra nước Việt Nam Dân chủ Cộng hòa vào ngày $2$ tháng $9$ năm $1945$. Năm đó thuộc thế kỉ nào và tính đến ngày $2$ tháng $9$ năm $2025$ là tròn bao nhiêu năm?
\choice
{Thế kỉ XIX; $70$ năm}
{\True Thế kỉ XX; $80$ năm}
{Thế kỉ XX; $75$ năm}
{Thế kỉ XXI; $80$ năm}
\loigiai{
Năm $1945$ thuộc thế kỉ XX (từ năm $1901$ đến $2000$).\\
Số năm tính đến ngày $2$ tháng $9$ năm $2025$ là:
$$2025 - 1945 = 80\text{ (năm)}.$$
\textbf{Đáp án của câu hỏi là B}
}
\end{ex}

\begin{ex} %[5D10V56]
Một vận động viên hoàn thành chặng đua xe đạp trong thời gian $2$ giờ $18$ phút. Thời gian đó viết dưới dạng số thập phân có đơn vị đo là giờ là
\choice
{$2{,}18\text{ giờ}$}
{\True $2{,}3\text{ giờ}$}
{$2{,}25\text{ giờ}$}
{$2{,}4\text{ giờ}$}
\loigiai{
Đổi: $18\text{ phút} = \dfrac{18}{60}\text{ giờ} = 0{,}3\text{ giờ}$.\\
Do đó $2\text{ giờ } 18\text{ phút} = 2 + 0{,}3 = 2{,}3\text{ giờ}$.\\
\textbf{Đáp án của câu hỏi là B}
}
\end{ex}

\begin{ex} %[5D10V56]
Một chiếc máy bay cất cánh lúc $21\text{ giờ } 45\text{ phút}$ ngày hôm trước và hạ cánh lúc $6\text{ giờ } 15\text{ phút}$ ngày hôm sau. Tổng thời gian chuyến bay đó kéo dài là
\choice
{$8\text{ giờ } 15\text{ phút}$}
{\True $8\text{ giờ } 30\text{ phút}$}
{$8\text{ giờ } 45\text{ phút}$}
{$9\text{ giờ } 30\text{ phút}$}
\loigiai{
Thời gian bay từ lúc cất cánh đến nửa đêm ($24\text{ giờ}$) là:
$$24\text{ giờ} - 21\text{ giờ } 45\text{ phút} = 2\text{ giờ } 15\text{ phút}.$$
Tổng thời gian bay là:
$$2\text{ giờ } 15\text{ phút} + 6\text{ giờ } 15\text{ phút} = 8\text{ giờ } 30\text{ phút}.$$
\textbf{Đáp án của câu hỏi là B}
}
\end{ex}

\begin{ex} %[5D10V56]
Năm $2024$ là năm nhuận. Hỏi từ ngày $1$ tháng $1$ năm $2024$ đến hết ngày $31$ tháng $12$ năm $2024$ có tất cả bao nhiêu tuần lễ và còn dư bao nhiêu ngày?
\choice
{$52\text{ tuần và dư } 1\text{ ngày}$}
{\True $52\text{ tuần và dư } 2\text{ ngày}$}
{$51\text{ tuần và dư } 3\text{ ngày}$}
{$52\text{ tuần không dư}$}
\loigiai{
Năm nhuận có $366$ ngày.\\
Thực hiện phép chia:
$$366 : 7 = 52\text{ (dư } 2\text{)}.$$
Vậy năm nhuận có $52$ tuần và dư $2$ ngày.\\
\textbf{Đáp án của câu hỏi là B}
}
\end{ex}

\begin{ex} %[5D10V56]
Một người thợ làm một chi tiết máy hết $45$ phút. Nếu người đó làm việc liên tục trong $6$ giờ thì làm được tất cả bao nhiêu chi tiết máy như vậy?
\choice
{$6\text{ chi tiết}$}
{\True $8\text{ chi tiết}$}
{$9\text{ chi tiết}$}
{$10\text{ chi tiết}$}
\loigiai{
Đổi: $6\text{ giờ} = 360\text{ phút}$.\\
Số chi tiết máy làm được là:
$$360 : 45 = 8\text{ (chi tiết)}.$$
\textbf{Đáp án của câu hỏi là B}
}
\end{ex}

\begin{ex} %[5D10V56]
Cho biết $\dfrac{3}{5}\text{ giờ} + 18\text{ phút} = \ldots\ldots\text{ phút}$. Số thích hợp điền vào chỗ chấm là
\choice
{$36$}
{\True $54$}
{$48$}
{$60$}
\loigiai{
Đổi $\dfrac{3}{5}\text{ giờ} = \dfrac{3}{5} \times 60\text{ phút} = 36\text{ phút}$.\\
Giá trị biểu thức là:
$$36 + 18 = 54\text{ (phút)}.$$
\textbf{Đáp án của câu hỏi là B}
}
\end{ex}

\begin{ex} %[5D10V56]
Một đồng hồ quả lắc cứ mỗi giờ lại chạy chậm $3$ giây. Hỏi sau một tuần lễ, chiếc đồng hồ đó chạy chậm tất cả bao nhiêu giây?
\choice
{$336\text{ giây}$}
{\True $504\text{ giây}$}
{$480\text{ giây}$}
{$560\text{ giây}$}
\loigiai{
Một tuần lễ có $7$ ngày, mỗi ngày có $24$ giờ.\\
Số giờ trong một tuần lễ là:
$$7 \times 24 = 168\text{ (giờ)}.$$
Số giây đồng hồ chạy chậm sau một tuần là:
$$168 \times 3 = 504\text{ (giây)}.$$
\textbf{Đáp án của câu hỏi là B}
}
\end{ex}

\begin{ex} %[5D10V56]
Khoảng thời gian từ $7\text{ giờ } 15\text{ phút}$ sáng đến $3\text{ giờ } 45\text{ phút}$ chiều cùng ngày là
\choice
{$8\text{ giờ}$}
{\True $8\text{ giờ } 30\text{ phút}$}
{$8\text{ giờ } 45\text{ phút}$}
{$9\text{ giờ } 30\text{ phút}$}
\loigiai{
Đổi: $3\text{ giờ } 45\text{ phút}$ chiều $= 15\text{ giờ } 45\text{ phút}$.\\
Khoảng thời gian đó là:
$$15\text{ giờ } 45\text{ phút} - 7\text{ giờ } 15\text{ phút} = 8\text{ giờ } 30\text{ phút}.$$
\textbf{Đáp án của câu hỏi là B}
}
\end{ex}

\begin{ex} %[5D10V56]
Số giây có trong $2{,}75\text{ giờ}$ là
\choice
{$9\,600\text{ giây}$}
{\True $9\,900\text{ giây}$}
{$10\,200\text{ giây}$}
{$8\,800\text{ giây}$}
\loigiai{
Một giờ có $3\,600\text{ giây}$.\\
Số giây trong $2{,}75$ giờ là:
$$2{,}75 \times 3\,600 = 9\,900\text{ (giây)}.$$
\textbf{Đáp án của câu hỏi là B}
}
\end{ex}

\begin{ex} %[5H9V56]
Biết ngày $15$ tháng $3$ của một năm là ngày thứ Bảy. Hỏi ngày $2$ tháng $4$ của năm đó là ngày thứ mấy?
\choice
{Thứ Ba}
{\True Thứ Tư}
{Thứ Năm}
{Thứ Sáu}
\loigiai{
Tháng 3 có $31$ ngày.\\
Số ngày từ ngày $15$ tháng $3$ đến ngày $2$ tháng $4$ là:
$$(31 - 15) + 2 = 16 + 2 = 18\text{ (ngày)}.$$
Ta có: $18 : 7 = 2$ (dư $4$).\\
Đếm thêm $4$ ngày kể từ thứ Bảy: Chủ nhật (1), Thứ Hai (2), Thứ Ba (3), Thứ Tư (4).\\
Vậy ngày $2$ tháng $4$ là ngày thứ Tư.\\
\textbf{Đáp án của câu hỏi là B}
}
\end{ex}

\begin{ex} %[5D10V56]
Một bộ phim hoạt hình dài $1$ giờ $45$ phút. Phim bắt đầu chiếu lúc $19\text{ giờ } 30\text{ phút}$ và có $15$ phút quảng cáo xen giữa. Hỏi bộ phim kết thúc vào lúc nào?
\choice
{$21\text{ giờ } 15\text{ phút}$}
{\True $21\text{ giờ } 30\text{ phút}$}
{$21\text{ giờ } 45\text{ phút}$}
{$22\text{ giờ}$}
\loigiai{
Tổng thời gian chiếu phim và quảng cáo là:
$$1\text{ giờ } 45\text{ phút} + 15\text{ phút} = 2\text{ giờ}.$$
Thời điểm bộ phim kết thúc là:
$$19\text{ giờ } 30\text{ phút} + 2\text{ giờ} = 21\text{ giờ } 30\text{ phút}.$$
\textbf{Đáp án của câu hỏi là B}
}
\end{ex}

\begin{ex} %[5D10V56]
Một vệ tinh nhân tạo bay một vòng quanh Trái Đất hết $1\text{ giờ } 36\text{ phút}$. Trong một ngày đêm ($24\text{ giờ}$), vệ tinh đó bay được tất cả bao nhiêu vòng?
\choice
{$12\text{ vòng}$}
{\True $15\text{ vòng}$}
{$16\text{ vòng}$}
{$18\text{ vòng}$}
\loigiai{
Đổi: $1\text{ giờ } 36\text{ phút} = 96\text{ phút}$.\\
Đổi: $24\text{ giờ} = 1\,440\text{ phút}$.\\
Số vòng vệ tinh bay được trong một ngày đêm là:
$$1\,440 : 96 = 15\text{ (vòng)}.$$
\textbf{Đáp án của câu hỏi là B}
}
\end{ex}

\textbf{Phần 2. Trắc nghiệm đúng sai}
\setcounterpageref{ex}{0}

\begin{ex} %[5D10V56]
Cho khoảng thời gian $T = 3\text{ giờ } 24\text{ phút}$. Xét tính đúng/sai của các nhận định sau:
\choiceTF
{\True $T$ đổi ra phút bằng $204$ phút}
{\True $T$ đổi ra giờ dưới dạng phân số tối giản là $\dfrac{17}{5}$ giờ}
{$T$ đổi ra số thập phân có đơn vị giờ là $3{,}24$ giờ}
{\True $T$ đổi ra giây bằng $12\,240$ giây}
\loigiai{
\begin{itemchoice}
\item Đúng. $3 \times 60 + 24 = 204$ phút.
\item Đúng. $3 + \dfrac{24}{60} = 3 + \dfrac{2}{5} = \dfrac{17}{5}$ giờ.
\item Sai. $3 + \dfrac{2}{5} = 3{,}4$ giờ (chứ không phải $3{,}24$ giờ).
\item Đúng. $204 \times 60 = 12\,240$ giây.
\end{itemchoice}
}
\end{ex}

\begin{ex} %[5D10V56]
Xét các khẳng định về năm và thế kỉ sau:
\choiceTF
{\True Thế kỉ XXI bắt đầu từ năm $2001$ và kết thúc vào năm $2100$}
{\True Các năm chia hết cho $4$ (như năm $2024, 2028$) thông thường là năm nhuận có $366$ ngày}
{Năm $1900$ là năm nhuận vì $1900$ chia hết cho $4$}
{\True Tháng Hai của năm nhuận có $29$ ngày, năm không nhuận có $28$ ngày}
\loigiai{
\begin{itemchoice}
\item Đúng. Thế kỉ XXI tính từ ngày 1/1/2001 đến hết ngày 31/12/2100.
\item Đúng. Các năm chia hết cho 4 (như 2024) có 366 ngày.
\item Sai. Đối với các năm tròn thế kỉ, chỉ năm chia hết cho 400 mới là năm nhuận (1900 không chia hết cho 400 nên không phải năm nhuận).
\item Đúng. Năm nhuận tháng Hai có 29 ngày, năm thường có 28 ngày.
\end{itemchoice}
}
\end{ex}

\textbf{Phần 3. Câu trả lời ngắn}
\setcounterpageref{ex}{0}

\begin{ex} %[5D10V56]
Một cuộc họp kéo dài từ $8\text{ giờ } 15\text{ phút}$ đến $11\text{ giờ } 45\text{ phút}$. Hỏi cuộc họp đó diễn ra trong bao nhiêu phút?
\shortans{210}
\loigiai{
Thời gian diễn ra cuộc họp là:
$$11\text{ giờ } 45\text{ phút} - 8\text{ giờ } 15\text{ phút} = 3\text{ giờ } 30\text{ phút}.$$
Đổi: $3\text{ giờ } 30\text{ phút} = 3 \times 60 + 30 = 210\text{ (phút)}$.
}
\end{ex}

\begin{ex} %[5D10V56]
Một người đạp xe $4$ vòng quanh một công viên hết $1$ giờ $20$ phút. Hỏi trung bình người đó đạp một vòng quanh công viên hết bao nhiêu phút?
\shortans{20}
\loigiai{
Đổi: $1\text{ giờ } 20\text{ phút} = 80\text{ phút}$.\\
Thời gian đạp một vòng là:
$$80 : 4 = 20\text{ (phút)}.$$
}
\end{ex}

\begin{ex} %[5D10V56]
Một vòi nước nhỏ giọt, cứ $12$ giây nhỏ được $1$ giọt nước. Biết cứ $20$ giọt nước thì được $1\text{ ml}$ nước. Hỏi trong $1$ giờ, vòi nước đó đã làm thất thoát bao nhiêu mi-li-lít nước?
\shortans{15}
\loigiai{
Đổi: $1\text{ giờ} = 3\,600\text{ giây}$.\\
Số giọt nước nhỏ trong $1$ giờ là:
$$3\,600 : 12 = 300\text{ (giọt)}.$$
Lượng nước thất thoát là:
$$300 : 20 = 15\text{ (ml)}.$$
}
\end{ex}

\begin{ex} %[5D10V56]
Một học sinh ngủ từ $22\text{ giờ}$ đêm hôm trước đến $6\text{ giờ } 30\text{ phút}$ sáng hôm sau. Hỏi bạn học sinh đó đã ngủ bao nhiêu phút?
\shortans{510}
\loigiai{
Thời gian ngủ từ $22\text{ giờ}$ đến nửa đêm ($24\text{ giờ}$) là $2\text{ giờ}$.\\
Thời gian ngủ từ nửa đêm đến $6\text{ giờ } 30\text{ phút}$ sáng là $6\text{ giờ } 30\text{ phút}$.\\
Tổng thời gian ngủ là:
$$2\text{ giờ} + 6\text{ giờ } 30\text{ phút} = 8\text{ giờ } 30\text{ phút}.$$
Đổi $8\text{ giờ } 30\text{ phút} = 8 \times 60 + 30 = 510\text{ (phút)}$.
}
\end{ex}

\begin{ex} %[5D10V56]
Một công nhân làm việc $5$ ngày mỗi tuần, mỗi ngày làm việc $8$ giờ. Hỏi sau $4$ tuần, người công nhân đó đã làm việc tất cả bao nhiêu giờ?
\shortans{160}
\loigiai{
Số ngày làm việc trong $4$ tuần là:
$$5 \times 4 = 20\text{ (ngày)}.$$
Tổng số giờ làm việc là:
$$20 \times 8 = 160\text{ (giờ)}.$$
}
\end{ex}

\begin{ex} %[5D10V56]
Chiến thắng Điện Biên Phủ lịch sử diễn ra vào năm $1954$. Năm đó thuộc thế kỉ thứ bao nhiêu?
\shortans{20}
\loigiai{
Năm $1954$ thuộc thế kỉ XX (thế kỉ thứ 20).
}
\end{ex}
